\problemname{Interactive problem}
\noindent

Innan Gösta dog av ålderdom kom vetenskapsmän på hur man kan göra människor odödliga. Hur detta gick till lämnas som en
kluring till läsaren. Efter $N$ år har gösta äntligen fylla sin åker med maskrosor (se \textit{Maskrosor}). Intressant nog
verkar de nya maskrosorna varje år vara smått annorlunda från tidigare år. De är så annorlunda att det finns $N$ olika sorters
maskrosor. Efter en del möda har Gösta lyckats plocka exakt $2$ stycken maskrosor av varje sort. Precis när han skulle börja
analysera de, så råkar han snubbla och blanda ihop alla maskrosorna. Katastrof!

Eftersom maskrosorna inte förändrats särskilt mycket visuellt kommer han behöva använda sin pollen-analys-maskin (PAM). 
Denna fungerar genom att låta alla maskrosor i den blomma, och kan sedan analysera pollenet i den för att räkna antalet
unika sorters maskrosor som finns i den. Detta beror på att mängden pollen kan variera väldigt mycket från maskros till maskros,
vilket leder till att den inte skilja på om det finns en eller två av samma sorts maskros i den. 

För att använda maskinen kan du göra två saker:
\begin{itemize}
  \item Lägga in en maskros i maskinen.
  \item Ta ut en maskros ur maskinen (som redan finns i den).
\end{itemize}

Gösta vill nu ha din hjälp att använda maskinen för att hitta vilka par av maskrosor som är från samma år.
Eftersom han är odödlig har han inte bråttom, men vill ändå bli färdig inom "rimlig" tid. Därför kan han acceptera
att du använder den som mest en miljon ($10^6$) gånger.

\section*{Indata}
Den första raden innehåller heltalet $N$ ($1 \leq N \leq 43000$), antalet år som Gösta plockade maskrosor från.

Därefter följer tre rader som vardera beskriver en rad av rutnätet. Den $i$:te raden innehåller heltalen
$c_1, c_2, \dots, c_N$ ($0 \leq c_j \leq c_N$), avståndet mellan svansen och ormens kroppssegment på rad $i$ och kolumn $j$.
Om $c_j$ är $0$ betyder det att Olle glömde vad $c_j$ är. Alla $c_j$ som inte är $0$ är parvis unika. Det är även garanterat att det
finns minst en giltig lösning.

\section*{Utdata}
Skriv ut $3$ rader med $N$ heltal vardera, som beskriver en giltig vinnande orm. Varje heltal mellan $1$ och $N \cdot 3$
måste förekomma exakt en gång. För varje $i$ mellan $1$ och $N \cdot 3 - 1$ måste segmentet med tal $i$ ligga ovanför, nedanför,
till vänster eller höger relativt segmentet med tal $i+1$.

\section*{Poängsättning}
Din lösning kommer att testas på en mängd testfallsgrupper.
För att få poäng för en grupp så måste du klara alla testfall i gruppen.

\noindent
\begin{tabular}{| l | l | p{12cm} |}
  \hline
  \textbf{Grupp} & \textbf{Poäng} & \textbf{Gränser} \\ \hline
  $1$    & $15$         & $N \leq 10$  \\ \hline
  $2$    & $25$         & $N \leq 40$ \\ \hline
  $3$    & $30$         & $N \leq 300$ \\ \hline
  $4$    & $30$         & Inga ytterligare begränsningar. \\ \hline
\end{tabular}

